% F085｜真实数学图；代表函数不构成抽象解的唯一性证明。
\begin{center}
\begin{tikzpicture}[x=1.8cm,y=1cm,>=stealth]
\draw[->] (0,0)--(5.6,0) node[right,font=\small] {输入};
\foreach \t in {1,2,3,4,5} {\draw (\t,.06)--(\t,-.06) node[below,font=\small] {$\t$};}
\fill[CExplanation] (2,0) circle(.045) (4,0) circle(.045);
\fill[orange!80!black] (3,0) circle(.045);
\draw[<->,thick,orange!80!black] (2,.55)--node[above,font=\small] {$y=1$}(3,.55);
\draw[<->,thick,orange!80!black] (3,.55)--node[above,font=\small] {$y=1$}(4,.55);
\node[font=\small,above] at (2,1.1) {$n=x-y=2$};
\node[font=\small,above] at (4,1.1) {$m=x+y=4$};
\node[font=\small,below] at (3,-.65) {中点 $x=3$，半差 $y=1$};
\end{tikzpicture}
\par\small 数轴上端点 $2,4$ 关于 $3$ 对称，所以直接代 $x=3,y=1$；无需猜测抽象函数的曲线。
\end{center}
